\documentclass[../../main.tex]{subfiles}

\graphicspath{{../../image/}{../../figure/}}

\begin{document}

\section{习题}\label{section:ch17-习题}

\begin{exercise}[$\star\star\star$]\label{exercise:ch17-01}
我们希望 GAN 损失函数［\eqref{align:ch17-GAN损失}］具有如下性质：在神经网络足够灵活的情况下，当生成分布与真实数据分布匹配时，就能获得驻点。在这个习题中，假设模型具有无限的灵活性，我们将通过在与生成器网络和判别器网络相对应的全部概率分布 $p_g(x)$ 和全部函数 $d(x)$ 的空间上对其进行优化来证明这一结果。具体来说，我们假设判别模型在一个内循环中被优化，从而产生用于生成式模型的有效外循环误差函数。首先证明在数据样本无限多的情况下，GAN 损失函数［\eqref{align:ch17-GAN损失}］可以重写为如下形式：
\begin{align}
E(p_g,d)=-\int p_{\mathrm{data}}(x)\ln d(x)\,\mathrm{d}x-\int p_g(x)\ln\big(1-d(x)\big)\,\mathrm{d}x,
\label{align:ch17-习题积分形式}
\end{align}
其中 $p_{\mathrm{data}}(x)$ 是真实数据点的固定分布。考虑对所有函数 $d(x)$ 进行变分优化（参见\hyperref[chapter:appB-变分法]{附录\ref*{chapter:appB-变分法}}），证明对于固定的生成器网络，使 $E$ 最小化的判别器函数 $d(x)$ 的解由下式给出：
\begin{align}
d^*(x)=\frac{p_{\mathrm{data}}(x)}{p_{\mathrm{data}}(x)+p_g(x)}.
\label{align:ch17-习题最优判别器}
\end{align}
进一步证明损失函数 $E$ 可以作为生成器网络 $p_g(x)$ 的函数，写成以下形式：
\begin{align}
C(p_g)={}&-\int p_{\mathrm{data}}(x)\ln\left\{\frac{p_{\mathrm{data}}(x)}{p_{\mathrm{data}}(x)+p_g(x)}\right\}\mathrm{d}x\nonumber\\
&-\int p_g(x)\ln\left\{\frac{p_g(x)}{p_{\mathrm{data}}(x)+p_g(x)}\right\}\mathrm{d}x,
\label{align:ch17-习题代价函数}
\end{align}
\eqref{align:ch17-习题代价函数}可以重写为以下形式：
\begin{align}
C(p_g)=-\ln(4)+\operatorname{KL}\left(p_{\mathrm{data}}\,\middle\|\,\frac{p_{\mathrm{data}}+p_g}{2}\right)+\operatorname{KL}\left(p_g\,\middle\|\,\frac{p_{\mathrm{data}}+p_g}{2}\right),
\label{align:ch17-习题KL形式}
\end{align}
其中 Kullback-Leibler 散度 $\operatorname{KL}(p\|q)$ 由\eqref{align:ch02-kl散度}定义（参见\cref{subsection:ch02-KL散度}）。最后，根据只有在 $p(x)=q(x)$ 对所有 $x$ 都成立时 $\operatorname{KL}(p\|q)$ 才等于零这一性质，证明 $C(p_g)$ 的最小值发生在 $p_g(x)=p_{\mathrm{data}}(x)$ 时。注意，\eqref{align:ch17-习题KL形式}中两个 Kullback-Leibler 散度的总和称为 $p_{\mathrm{data}}$ 和 $p_g$ 之间的 Jensen-Shannon 散度。和 Kullback-Leibler 散度一样，Jensen-Shannon 散度也是一个非负的量，当且仅当两个分布相等时它会消失，但不同于 KL 散度，它对于两个分布是对称的。
\end{exercise}

\begin{remark}（重构新增）\textbf{原文问题说明。}
本题的\eqref{align:ch17-习题代价函数}与\eqref{align:ch17-习题KL形式}在符号上不相容。按\eqref{align:ch17-习题积分形式}的定义，$E(p_g,d^*)$ 的值为 $\ln4-\operatorname{JSD}$，而\eqref{align:ch17-习题KL形式}的右端为 $-\ln4+\operatorname{JSD}$，两者互为相反数。这一点可以用一个具体的例子核对：取 $p_{\mathrm{data}}=(1/2,1/2)$、$p_g=(1,0)$ 的两点分布，则 $d^*=(1/3,1)$，直接代入\eqref{align:ch17-习题代价函数}得 $E=\tfrac12\ln3+\ln\tfrac32\approx0.955$，而 $-\ln4+\operatorname{JSD}\approx-0.955$，二者确实相差一个符号。

最小修正如下（保留原题的条件与结论，只修正符号约定）：把\eqref{align:ch17-习题代价函数}改写为“$C(p_g)\triangleq-E(p_g,d^*)$”，则\eqref{align:ch17-习题KL形式}成立，且“$C$ 的最小值在 $p_g=p_{\mathrm{data}}$ 处取到”成立；或者保留\eqref{align:ch17-习题代价函数}不动，把\eqref{align:ch17-习题KL形式}写成 $C(p_g)=\ln4-\operatorname{JSD}(p_{\mathrm{data}}\|p_g)$，并把结论改为“$C$ 的最大值在 $p_g=p_{\mathrm{data}}$ 处取到，等价地 JSD 取最小值零”。本工程的正文采用后一种说法，见\cref{proposition:ch17-最优判别器}。
\end{remark}

\begin{exercise}[$\star\star\star$]\label{exercise:ch17-02}
在这个习题中，我们将探讨 GAN 训练的对抗性可能会引起的问题。考虑一个定义在两个参数 $a$ 和 $b$ 上的成本函数 $E(a,b)=ab$，证明点 $a=0$，$b=0$ 是该成本函数的一个驻点。考虑沿着直线 $b=a$ 和 $b=-a$ 的二阶导数，其中 $a$ 和 $b$ 分别类比于生成网络和判别网络的参数。证明点 $a=0$，$b=0$ 也是一个鞍点。假设我们采取无穷小步骤来优化这个误差函数，使得变量成为连续时间 $a(t)$ 和 $b(t)$ 的函数，在此过程中生成器网络的参数 $a(t)$ 被更新以增加 $E(a,b)$，参数 $b(t)$ 也被更新以减少 $E(a,b)$。证明参数的演变受以下方程的控制：
\begin{align}
\frac{\mathrm{d}a}{\mathrm{d}t}=\eta\frac{\partial E}{\partial a},\qquad
\frac{\mathrm{d}b}{\mathrm{d}t}=-\eta\frac{\partial E}{\partial b}.
\label{align:ch17-鞍点动力学}
\end{align}
然后证明 $a(t)$ 满足二阶微分方程：
\begin{align}
\frac{\mathrm{d}^2a}{\mathrm{d}t^2}=-\eta^2a(t).
\label{align:ch17-二阶方程}
\end{align}
验证以上二阶微分方程的解为
\begin{align}
a(t)=C\cos(\eta t)+D\sin(\eta t),
\label{align:ch17-二阶方程解}
\end{align}
其中 $C$ 和 $D$ 是任意常数。如果系统在 $t=0$ 时初始化为 $a=1$，$b=0$，找出 $C$ 和 $D$ 的值，并由此证明得到的 $a(t)$ 和 $b(t)$ 的值会在以原点为中心的 $a,b$ 空间中描绘出一个单位半径的圆，并且它们因此永远不会收敛到鞍点。
\end{exercise}

\begin{exercise}[$\star\star$]\label{exercise:ch17-03}
考虑一个 GAN 模型，其训练集由相等数量的猫和狗的图像构成，并且生成器网络已经学会了产生高质量的狗的图像。证明当我们向判别器网络（训练目标是输出一个图像为真的概率）输入一张狗的图像时，最优输出是 $1/3$。
\end{exercise}

\end{document}
